\section{Results}
\label{sec:results}

\subsection{Main comparison results on public ultrasound datasets}
Table~\ref{tab:main_results} summarizes the main comparison results, reported as mean $\pm$ standard deviation across TN5000 seeds or BUSI/AUL training folds evaluated on the fixed held-out test sets. The evidence should be read carefully. On TN5000, ConvNeXt-Tiny is slightly stronger than the proposed model at the selected operating point. On BUSI and AUL, however, UL-TransXNet provides a clear gain over the strongest non-\textit{Ours} baseline. Averaged over the three dataset-specific evaluations, the proposed model gives the strongest overall performance, especially in balanced accuracy.

\begin{table*}[t]
    \centering
    \small
    \caption{Main comparison results on the public ultrasound datasets. For each dataset, the table reports the proposed model and the strongest non-\textit{Ours} baseline ranked by mean test AUC. F1 denotes macro F1.}
    \label{tab:main_results}
    \resizebox{\textwidth}{!}{%
    \begin{tabular}{llcccc}
        \toprule
        Dataset & Method & AUC & BalAcc & F1-macro & Acc \\
        \midrule
        TN5000 & UL-TransXNet & 0.9483 $\pm$ 0.0006 & 0.8726 $\pm$ 0.0046 & 0.8503 $\pm$ 0.0123 & 0.8790 $\pm$ 0.0140 \\
        TN5000 & ConvNeXt-Tiny & \textbf{0.9508 $\pm$ 0.0020} & \textbf{0.8794 $\pm$ 0.0023} & \textbf{0.8731 $\pm$ 0.0066} & \textbf{0.8957 $\pm$ 0.0065} \\
        \midrule
        BUSI & UL-TransXNet & \textbf{0.8998 $\pm$ 0.0116} & \textbf{0.7961 $\pm$ 0.0104} & \textbf{0.8018 $\pm$ 0.0119} & \textbf{0.8140 $\pm$ 0.0130} \\
        BUSI & Swin-T & 0.8190 $\pm$ 0.0091 & 0.7137 $\pm$ 0.0095 & 0.6896 $\pm$ 0.0300 & 0.7194 $\pm$ 0.0445 \\
        \midrule
        AUL & UL-TransXNet & \textbf{0.8767 $\pm$ 0.0009} & \textbf{0.8200 $\pm$ 0.0487} & 0.7462 $\pm$ 0.0405 & 0.7685 $\pm$ 0.0418 \\
        AUL & Swin-T & 0.8319 $\pm$ 0.0117 & 0.7783 $\pm$ 0.0325 & \textbf{0.7589 $\pm$ 0.0243} & \textbf{0.7717 $\pm$ 0.0211} \\
        \bottomrule
    \end{tabular}%
    }
\end{table*}

These results delimit the central claim. The proposed model should not be described as uniformly better than every baseline on every dataset. Its advantage is a multi-dataset in-domain pattern: it is competitive on TN5000 and substantially stronger on BUSI and AUL under dataset-specific training. The mean AUC over TN5000, BUSI, and AUL is 0.9083 for UL-TransXNet, compared with 0.8672 for the strongest-per-dataset non-\textit{Ours} rows listed in Table~\ref{tab:main_results}. The corresponding mean balanced accuracy is 0.8296 versus 0.7905. This supports the claim that the architecture is a strong compact multi-dataset ultrasound classifier, not a claim of uniform dominance on the primary thyroid benchmark or external cross-organ generalization.

AUL illustrates the distinction between threshold-free ranking and thresholded operating-point metrics. UL-TransXNet has higher AUC and balanced accuracy than Swin-T, whereas Swin-T has slightly higher macro F1 and raw accuracy at its validation-selected threshold. We therefore interpret the AUL result as better ranking and class-balanced operating behavior, not as dominance over every thresholded metric.

\begin{table}[t]
    \centering
    \small
    \caption{Claim boundary implied by the current evidence.}
    \label{tab:claim_boundary}
    \begin{adjustbox}{max width=\columnwidth}
    \begin{tabular}{lll}
        \toprule
        Claim & Supported by & Not supported by this study \\
        \midrule
        Multi-dataset in-domain competitiveness & Dataset-specific TN5000/BUSI/AUL tests & External cross-organ transfer \\
        Compact ROI classification & Params/MACs and ROI benchmark results & Smallest or CUDA-fastest model \\
        Validation-audited model selection & Mean validation AUC protocol and checksums & Test-set architecture search \\
        Automatic ROI workflow feasibility & TN5000 detector check; BUSI/AUL auto vs full & Novel lesion detector \\
        Mobile feasibility prototype & Android ONNX Runtime measurement & Clinical deployment validation \\
        \bottomrule
    \end{tabular}
    \end{adjustbox}
\end{table}

Table~\ref{tab:paired_statistical_tests} adds an offline paired statistical check over the same repeated runs used by Table~\ref{tab:main_results}. On TN5000, the proposed model is slightly below ConvNeXt-Tiny. On BUSI and AUL, the paired differences are consistently positive for both AUC and balanced accuracy. Because the available repeated units are only three TN5000 seeds and five BUSI/AUL folds, these tests are used as stability evidence rather than as large-sample clinical significance claims.

Table~\ref{tab:original_transxnet_delta} directly compares the original TransXNet baseline with UL-TransXNet under the same dataset-specific reporting protocol. This comparison is important because the method is a TransXNet-family adaptation rather than a detached backbone. The final model is not a TN5000 AUC improvement over the original TransXNet, but it slightly improves the TN5000 validation-selected operating point and gives much larger gains on BUSI and AUL.

\begin{table}[t]
    \centering
    \small
    \caption{Original TransXNet versus UL-TransXNet under the same dataset-specific reporting protocol. $\Delta$ denotes UL-TransXNet minus TransXNet.}
    \label{tab:original_transxnet_delta}
    \begin{adjustbox}{max width=\columnwidth}
    \begin{tabular}{lcccccc}
        \toprule
        Dataset & \multicolumn{2}{c}{TransXNet} & \multicolumn{2}{c}{UL-TransXNet} & \multicolumn{2}{c}{$\Delta$} \\
        \cmidrule(lr){2-3}\cmidrule(lr){4-5}\cmidrule(lr){6-7}
        & AUC & BalAcc & AUC & BalAcc & AUC & BalAcc \\
        \midrule
        TN5000 & 0.9495 & 0.8697 & 0.9483 & 0.8726 & -0.0012 & 0.0028 \\
        BUSI & 0.7227 & 0.6560 & 0.8998 & 0.7961 & 0.1770 & 0.1401 \\
        AUL & 0.6582 & 0.6263 & 0.8767 & 0.8200 & 0.2186 & 0.1937 \\
        \bottomrule
    \end{tabular}
    \end{adjustbox}
\end{table}

\begin{table}[t]
    \centering
    \small
    \caption{Offline paired statistical comparison between UL-TransXNet and the strongest non-\textit{Ours} baseline in Table~\ref{tab:main_results}. $\Delta$ denotes UL-TransXNet minus baseline. Confidence intervals are paired bootstrap 95\% intervals over seeds or folds; $p$ is an exact two-sided paired sign-flip test. Because only three TN5000 seeds and five BUSI/AUL folds are available, these tests are intended as repeated-run stability checks rather than large-sample clinical significance tests.}
    \label{tab:paired_statistical_tests}
    \begin{adjustbox}{max width=\columnwidth}
    \begin{tabular}{lllcccc}
        \toprule
        Dataset & Baseline & Metric & $n$ & $\Delta$ & 95\% CI & $p$ \\
        \midrule
        TN5000 & ConvNeXt-Tiny & AUC & 3 & -0.0025 & [-0.0040, -0.0006] & 0.2500 \\
        TN5000 & ConvNeXt-Tiny & BalAcc & 3 & -0.0068 & [-0.0104, -0.0022] & 0.2500 \\
        \midrule
        BUSI & Swin-T & AUC & 5 & +0.0807 & [0.0697, 0.0918] & 0.0625 \\
        BUSI & Swin-T & BalAcc & 5 & +0.0823 & [0.0647, 0.0961] & 0.0625 \\
        \midrule
        AUL & Swin-T & AUC & 5 & +0.0449 & [0.0375, 0.0574] & 0.0625 \\
        AUL & Swin-T & BalAcc & 5 & +0.0417 & [0.0160, 0.0674] & 0.0625 \\
        \bottomrule
    \end{tabular}
    \end{adjustbox}
\end{table}

Table~\ref{tab:complexity_performance} gives the corresponding complexity--performance view for the full comparison pool. This table makes the trade-off more explicit than the accuracy-only table: UL-TransXNet is parameter/MAC-efficient, but it is not the smallest model and is not CUDA-latency optimal. It achieves the highest mean AUC and mean balanced accuracy with 14.37M parameters and 2.36 GMACs, while its local CUDA batch-1 latency is 30.23 ms, slower than several GPU-optimized baselines. In contrast, the strongest TN5000 baselines ConvNeXt-Tiny and Swin-T require 28.21M/5.82G and 27.91M/7.11G, respectively. The input column is included because EfficientFormer-L1 is evaluated at its native 224 input while the other models use the 256 ROI input used by the proposed model.

\begin{table*}[t]
    \centering
    \small
    \caption{Complexity--performance trade-off of the comparison pool. CUDA latency is the mean model-only batch-1 latency measured on the local workstation; AUC values are mean test AUCs across TN5000 seeds or BUSI/AUL folds. Mean AUC and Mean BalAcc are averaged over the three datasets.}
    \label{tab:complexity_performance}
    \resizebox{\textwidth}{!}{%
    \begin{tabular}{lllccccccccc}
        \toprule
        Method & Type & Input & Params (M)$\downarrow$ & MACs (G)$\downarrow$ & CUDA ms$\downarrow$ & CUDA FPS$\uparrow$ & \multicolumn{3}{c}{AUC$\uparrow$} & Mean AUC$\uparrow$ & Mean BalAcc$\uparrow$ \\
        \cmidrule(lr){8-10}
        & & & & & & & TN5000 & BUSI & AUL & & \\
        \midrule
        MobileNetV3-Large & Mobile CNN & $256 \times 256$ & 4.86 & 0.28 & 3.17 & 315.1 & 0.8282 & 0.7947 & 0.7357 & 0.7862 & 0.7183 \\
        EfficientNet-B0 & Efficient CNN & $256 \times 256$ & 4.66 & 0.50 & 6.53 & 153.1 & 0.7856 & 0.7659 & 0.6099 & 0.7205 & 0.6682 \\
        EfficientFormer-L1 & Hybrid & $224 \times 224$ & 11.62 & 1.31 & 3.41 & 293.0 & 0.9249 & 0.7917 & 0.7196 & 0.8121 & 0.7330 \\
        RepViT-M1.1 & Mobile CNN & $256 \times 256$ & 8.04 & 1.80 & 5.19 & 192.5 & 0.8403 & 0.8028 & 0.6797 & 0.7743 & 0.7020 \\
        \textbf{UL-TransXNet} & \textbf{Ours} & $256 \times 256$ & 14.37 & 2.36 & 30.23 & 33.1 & 0.9483 & \textbf{0.8998} & \textbf{0.8767} & \textbf{0.9083} & \textbf{0.8296} \\
        DenseNet121 & CNN & $256 \times 256$ & 7.48 & 3.70 & 7.19 & 139.1 & 0.9360 & 0.8154 & 0.6729 & 0.8081 & 0.7401 \\
        ResNet50 & CNN & $256 \times 256$ & 24.56 & 5.40 & 3.25 & 307.5 & 0.6531 & 0.7519 & 0.5481 & 0.6510 & 0.6190 \\
        ConvNeXt-Tiny & ConvNet & $256 \times 256$ & 28.21 & 5.82 & 3.30 & 303.3 & \textbf{0.9508} & 0.8050 & 0.8274 & 0.8611 & 0.7651 \\
        Swin-T & Transformer & $256 \times 256$ & 27.91 & 7.11 & 6.32 & 158.3 & 0.9497 & 0.8190 & 0.8319 & 0.8669 & 0.7901 \\
        \bottomrule
    \end{tabular}%
    }
\end{table*}

\subsection{TN5000 progressive-family ablation}
Table~\ref{tab:tn5000_structure_ablation} reports the unified-protocol TN5000 structure ablation. The strongest TN5000 structure is \texttt{TransXNet-GG}, which achieves the highest AUC, balanced accuracy, macro F1, and accuracy in this ablation group. The \texttt{TransXNet-GGG-MCA} ablation row, which corresponds to UL-TransXNet, is close but not best on TN5000 alone. This distinction is important because it prevents the paper from over-attributing all gains to the final GGG-MCA setting.

\begin{table}[t]
    \centering
    \small
    \caption{Unified-protocol structure ablation on TN5000. Mean and standard deviation are computed over three seeds.}
    \label{tab:tn5000_structure_ablation}
    \begin{adjustbox}{max width=\columnwidth}
    \begin{tabular}{lcccc}
        \toprule
        Variant & AUC & BalAcc & F1-macro & Acc \\
        \midrule
        TransXNet & 0.9495 $\pm$ 0.0020 & 0.8697 $\pm$ 0.0100 & 0.8366 $\pm$ 0.0222 & 0.8640 $\pm$ 0.0230 \\
        TransXNet-G & 0.9481 $\pm$ 0.0013 & 0.8745 $\pm$ 0.0074 & 0.8395 $\pm$ 0.0220 & 0.8660 $\pm$ 0.0233 \\
        TransXNet-GG & \textbf{0.9517 $\pm$ 0.0011} & \textbf{0.8784 $\pm$ 0.0076} & \textbf{0.8552 $\pm$ 0.0173} & \textbf{0.8827 $\pm$ 0.0170} \\
        TransXNet-GGG-noMCA & 0.9503 $\pm$ 0.0019 & 0.8660 $\pm$ 0.0069 & 0.8325 $\pm$ 0.0013 & 0.8610 $\pm$ 0.0016 \\
        TransXNet-GGG-MCA & 0.9483 $\pm$ 0.0006 & 0.8726 $\pm$ 0.0046 & 0.8503 $\pm$ 0.0123 & 0.8790 $\pm$ 0.0140 \\
        \bottomrule
    \end{tabular}
    \end{adjustbox}
\end{table}

The structure ablation is a progressive-family analysis rather than a full factorial proof. It shows that the TransXNet-family modifications are useful but metric-dependent. The G and GG variants improve thresholded operating-point behavior relative to the baseline. The GGG-MCA configuration is not the TN5000 AUC winner, but it remains competitive and corresponds to the best multi-dataset in-domain pattern in this result set. Table~\ref{tab:validation_selection_audit} audits the model-family choice using validation AUC from the completed runs. The selection rule is to choose the TransXNet-family row with the highest mean validation AUC over TN5000, BUSI, and AUL under the dataset-specific training protocols; the source-controlled model-selection protocol records the contributing CSV files and checksums. UL-TransXNet gives the highest mean validation AUC across TN5000, BUSI, and AUL, while \texttt{GGG-noMCA} remains the strongest TN5000 validation-AUC row. Table~\ref{tab:gg_vs_final_selection} then shows the corresponding test-set behavior: \texttt{TransXNet-GG} is the strongest TN5000-only structure, whereas UL-TransXNet is substantially stronger on BUSI and AUL and therefore gives the better three-dataset average.

\begin{table}[t]
    \centering
    \small
    \caption{Validation-performance audit for the TransXNet-family selection. Values are averages over three TN5000 seeds or five BUSI/AUL training folds. The selected row is the final UL-TransXNet configuration used in the main manuscript tables.}
    \label{tab:validation_selection_audit}
    \begin{adjustbox}{max width=\columnwidth}
    \begin{tabular}{lccccc}
        \toprule
        Variant & TN5000 val AUC & BUSI val AUC & AUL val AUC & Mean val AUC & Selected \\
        \midrule
        TransXNet & 0.9577 & 0.7796 & 0.7464 & 0.8279 & no \\
        TransXNet-GG & 0.9632 & 0.8011 & 0.7390 & 0.8345 & no \\
        GGG-noMCA & \textbf{0.9700} & 0.8957 & 0.8798 & 0.9152 & no \\
        UL-TransXNet & 0.9638 & \textbf{0.9199} & \textbf{0.8827} & \textbf{0.9222} & yes \\
        \bottomrule
    \end{tabular}
    \end{adjustbox}
\end{table}

\begin{table}[t]
    \centering
    \small
    \caption{TN5000-only optimum versus the final multi-dataset configuration. The final reported model is retained for three-dataset behavior, not for the single best TN5000 ablation row.}
    \label{tab:gg_vs_final_selection}
    \begin{adjustbox}{max width=\columnwidth}
    \begin{tabular}{lcccc}
        \toprule
        Dataset & \multicolumn{2}{c}{TransXNet-GG} & \multicolumn{2}{c}{UL-TransXNet} \\
        \cmidrule(lr){2-3}\cmidrule(lr){4-5}
        & AUC & BalAcc & AUC & BalAcc \\
        \midrule
        TN5000 & \textbf{0.9517} & \textbf{0.8784} & 0.9483 & 0.8726 \\
        BUSI & 0.7707 & 0.7000 & \textbf{0.8998} & \textbf{0.7961} \\
        AUL & 0.6972 & 0.6550 & \textbf{0.8767} & \textbf{0.8200} \\
        \midrule
        Mean & 0.8065 & 0.7445 & \textbf{0.9083} & \textbf{0.8296} \\
        \bottomrule
    \end{tabular}
    \end{adjustbox}
\end{table}

\subsection{TN5000 core-module ablation}
Table~\ref{tab:tn5000_module_ablation} isolates the core-module path under the same TN5000 protocol. Differential attention provides the clearest single-module operating-point improvement: \texttt{GGG-core+DA} improves balanced accuracy, macro F1, and accuracy over \texttt{GGG-core}. MUDD alone gives a smaller gain in thresholded metrics and does not improve AUC in this setting. Adding MCA to the full GGG configuration lowers TN5000 AUC slightly relative to the no-MCA full model, but improves balanced accuracy, macro F1, and accuracy. This supports a restrained interpretation of MCA: it is useful for decision-level robustness and multi-dataset behavior, not a universal AUC booster.

\begin{table}[t]
    \centering
    \small
    \caption{Unified-protocol core-module ablation on TN5000. Mean and standard deviation are computed over three seeds.}
    \label{tab:tn5000_module_ablation}
    \begin{adjustbox}{max width=\columnwidth}
    \begin{tabular}{lcccc}
        \toprule
        Variant & AUC & BalAcc & F1-macro & Acc \\
        \midrule
        GGG-core & 0.9496 $\pm$ 0.0003 & 0.8704 $\pm$ 0.0083 & 0.8324 $\pm$ 0.0136 & 0.8593 $\pm$ 0.0133 \\
        GGG-core+MUDD & 0.9470 $\pm$ 0.0006 & 0.8712 $\pm$ 0.0045 & 0.8406 $\pm$ 0.0097 & 0.8687 $\pm$ 0.0111 \\
        GGG-core+DA & 0.9489 $\pm$ 0.0016 & \textbf{0.8751 $\pm$ 0.0044} & 0.8490 $\pm$ 0.0100 & 0.8770 $\pm$ 0.0099 \\
        GGG-core+MUDD+DA-noMCA & \textbf{0.9503 $\pm$ 0.0019} & 0.8660 $\pm$ 0.0069 & 0.8325 $\pm$ 0.0013 & 0.8610 $\pm$ 0.0016 \\
        GGG-core+MUDD+DA+MCA & 0.9483 $\pm$ 0.0006 & 0.8726 $\pm$ 0.0046 & \textbf{0.8503 $\pm$ 0.0123} & \textbf{0.8790 $\pm$ 0.0140} \\
        \bottomrule
    \end{tabular}
    \end{adjustbox}
\end{table}

\subsection{MCA effect across datasets}
Table~\ref{tab:mca_effect} compares the final GGG model with and without MCA across the three main datasets. The result explains why the final manuscript should keep MCA in the model family but avoid overclaiming it. On TN5000, MCA improves balanced accuracy, macro F1, and accuracy while slightly reducing AUC. On BUSI, MCA provides a large gain across all reported metrics. On AUL, MCA gives a small AUC and balanced-accuracy gain with nearly unchanged macro F1 and accuracy. Thus, MCA is best described as a module that improves multi-dataset operating behavior, especially on BUSI, rather than as a universally dominant change.

\begin{table*}[t]
    \centering
    \small
    \caption{Effect of enabling MCA in the final GGG configuration. Mean and standard deviation are computed over TN5000 seeds or BUSI/AUL folds.}
    \label{tab:mca_effect}
    \resizebox{\textwidth}{!}{%
    \begin{tabular}{llcccc}
        \toprule
        Dataset & Variant & AUC & BalAcc & F1-macro & Acc \\
        \midrule
        TN5000 & GGG-noMCA & \textbf{0.9503 $\pm$ 0.0019} & 0.8660 $\pm$ 0.0069 & 0.8325 $\pm$ 0.0013 & 0.8610 $\pm$ 0.0016 \\
        TN5000 & GGG-MCA & 0.9483 $\pm$ 0.0006 & \textbf{0.8726 $\pm$ 0.0046} & \textbf{0.8503 $\pm$ 0.0123} & \textbf{0.8790 $\pm$ 0.0140} \\
        \midrule
        BUSI & GGG-noMCA & 0.8588 $\pm$ 0.0137 & 0.7551 $\pm$ 0.0295 & 0.7526 $\pm$ 0.0300 & 0.7628 $\pm$ 0.0288 \\
        BUSI & GGG-MCA & \textbf{0.8998 $\pm$ 0.0116} & \textbf{0.7961 $\pm$ 0.0104} & \textbf{0.8018 $\pm$ 0.0119} & \textbf{0.8140 $\pm$ 0.0130} \\
        \midrule
        AUL & GGG-noMCA & 0.8720 $\pm$ 0.0046 & 0.8181 $\pm$ 0.0338 & 0.7462 $\pm$ 0.0283 & 0.7685 $\pm$ 0.0321 \\
        AUL & GGG-MCA & \textbf{0.8767 $\pm$ 0.0009} & \textbf{0.8200 $\pm$ 0.0487} & 0.7462 $\pm$ 0.0405 & 0.7685 $\pm$ 0.0418 \\
        \bottomrule
    \end{tabular}%
    }
\end{table*}

\subsection{Automatic ROI localization and closed-loop classification}
The main ROI protocol isolates classifier behavior using annotation-derived crops. Figure~\ref{fig:auto_roi_workflow}, Table~\ref{tab:tn5000_auto_roi_summary}, Table~\ref{tab:auto_roi_detector_quality}, and Table~\ref{tab:auto_roi_closed_loop} evaluate a more complete inference chain: an automatic detector first predicts the lesion box, an expanded ROI is constructed from that prediction, and the same trained UL-TransXNet classifiers are then applied without retraining. The TN5000 automatic-ROI sanity check confirms that the primary benchmark can also be connected to a detector front end. For BUSI and AUL, the detector does not miss any test image in this evaluation, but localization quality differs by dataset. BUSI reaches a mean IoU of 0.7365, whereas AUL reaches 0.5755 and has a lower high-IoU recall.

\begin{figure*}[t]
    \centering
    \includegraphics[width=\textwidth]{fig_auto_roi_workflow.pdf}
    \caption{Automatic ROI inference workflow for the BUSI/AUL closed-loop evaluation. Green boxes indicate reference annotations used only for evaluation and visualization; red boxes are detector-predicted lesion boxes; blue dashed boxes are the expanded ROIs used as classifier inputs.}
    \label{fig:auto_roi_workflow}
\end{figure*}

\begin{table}[t]
    \centering
    \small
    \caption{TN5000 automatic ROI sanity check using existing detector and classifier outputs. Detector values summarize three YOLO11n detector seeds; classification values summarize automatic-ROI closed-loop evaluation across the three detector seeds. No classifier retraining is involved.}
    \label{tab:tn5000_auto_roi_summary}
    \begin{adjustbox}{max width=\columnwidth}
    \begin{tabular}{lcc}
        \toprule
        Block & Metric & Value \\
        \midrule
        Detector & mAP50 & 0.9520 $\pm$ 0.0090 \\
        Detector & mean IoU & 0.7855 $\pm$ 0.0033 \\
        Detector & R@0.75 & 0.7880 $\pm$ 0.0122 \\
        Detector & no-detection rate & 0.0000 $\pm$ 0.0000 \\
        \midrule
        Auto ROI classifier & AUC & 0.9437 $\pm$ 0.0042 \\
        Auto ROI classifier & BalAcc & 0.8686 $\pm$ 0.0056 \\
        Auto ROI classifier & F1-macro & 0.8519 $\pm$ 0.0116 \\
        Auto ROI classifier & Acc & 0.8820 $\pm$ 0.0115 \\
        \bottomrule
    \end{tabular}
    \end{adjustbox}
\end{table}

\begin{table}[t]
    \centering
    \small
    \caption{Detector localization quality used for automatic ROI construction. Values are mean $\pm$ standard deviation across five detector folds.}
    \label{tab:auto_roi_detector_quality}
    \begin{adjustbox}{max width=\columnwidth}
    \begin{tabular}{lcccc}
        \toprule
        Dataset & Mean IoU & R@0.50 & R@0.75 & No det. \\
        \midrule
        BUSI & 0.7365 $\pm$ 0.0241 & 0.8403 $\pm$ 0.0141 & 0.6822 $\pm$ 0.0548 & 0.0000 \\
        AUL & 0.5755 $\pm$ 0.0221 & 0.6740 $\pm$ 0.0419 & 0.4331 $\pm$ 0.0201 & 0.0000 \\
        \bottomrule
    \end{tabular}
    \end{adjustbox}
\end{table}

\begin{table*}[t]
    \centering
    \small
    \caption{Closed-loop classification with oracle ROI, automatic detector ROI, and full-image input. Values are mean $\pm$ standard deviation across five folds. Bracketed intervals are fold-bootstrap 95\% confidence intervals for AUC and balanced accuracy.}
    \label{tab:auto_roi_closed_loop}
    \resizebox{\textwidth}{!}{%
    \begin{tabular}{llcccc}
        \toprule
        Dataset & Input & AUC & BalAcc & F1-macro & Acc \\
        \midrule
        BUSI & oracle & 0.8998 $\pm$ 0.0116 [0.8880, 0.9067] & 0.7961 $\pm$ 0.0104 [0.7865, 0.8045] & 0.8018 $\pm$ 0.0119 & 0.8140 $\pm$ 0.0130 \\
        BUSI & auto & 0.8725 $\pm$ 0.0158 [0.8564, 0.8845] & 0.7839 $\pm$ 0.0179 [0.7704, 0.8010] & 0.7893 $\pm$ 0.0176 & 0.8031 $\pm$ 0.0167 \\
        BUSI & full & 0.7737 $\pm$ 0.0044 [0.7701, 0.7777] & 0.6982 $\pm$ 0.0237 [0.6762, 0.7178] & 0.6915 $\pm$ 0.0378 & 0.7023 $\pm$ 0.0418 \\
        \midrule
        AUL & oracle & 0.8767 $\pm$ 0.0009 [0.8759, 0.8776] & 0.8200 $\pm$ 0.0487 [0.7755, 0.8599] & 0.7462 $\pm$ 0.0405 & 0.7685 $\pm$ 0.0418 \\
        AUL & auto & 0.8507 $\pm$ 0.0222 [0.8320, 0.8708] & 0.7593 $\pm$ 0.0288 [0.7340, 0.7847] & 0.6908 $\pm$ 0.0561 & 0.7165 $\pm$ 0.0713 \\
        AUL & full & 0.7441 $\pm$ 0.0111 [0.7344, 0.7538] & 0.6805 $\pm$ 0.0214 [0.6623, 0.6999] & 0.6787 $\pm$ 0.0203 & 0.7512 $\pm$ 0.0190 \\
        \bottomrule
    \end{tabular}%
    }
\end{table*}

Automatic ROI classification is consistently stronger than full-image classification. On BUSI, the automatic ROI improves AUC by 0.0988 and balanced accuracy by 0.0857 over the full-image input. On AUL, the corresponding gains are 0.1066 AUC and 0.0789 balanced accuracy. The automatic ROI remains below oracle ROI, especially on AUL, which is consistent with the lower localization quality in Table~\ref{tab:auto_roi_detector_quality}. These results support the practical use of an automatic ROI front end, while keeping the detector itself outside the main architectural contribution.

\subsection{Qualitative Grad-CAM sanity check}
Figure~\ref{fig:gradcam_sanity} provides a qualitative Grad-CAM sanity check using one trained UL-TransXNet checkpoint per dataset. The maps are generated from late spatial backbone blocks selected for readable ROI-level maps. The activations broadly fall on lesion tissue and adjacent ultrasound context inside the ROI. This visualization is used only as a qualitative sanity check that the classifier is not driven solely by image borders or unrelated background. It is not used as quantitative evidence of causal attention or as a substitute for the ablation and closed-loop ROI experiments above.

\begin{figure}[t]
    \centering
    \includegraphics[width=0.88\textwidth]{fig_gradcam_sanity.pdf}
    \caption{Grad-CAM sanity-check visualization for representative correctly classified ROI inputs. Columns correspond to datasets; the upper row shows the ROI input and the lower row shows the Grad-CAM overlay. The maps are generated from a single UL-TransXNet checkpoint for each dataset and are intended as qualitative localization checks only.}
    \label{fig:gradcam_sanity}
\end{figure}

\subsection{Uncertainty and calibration analysis}
Table~\ref{tab:calibration_uncertainty} reports bootstrap confidence intervals and calibration metrics for UL-TransXNet using the same seed/fold-level reporting protocol as the main tables. The confidence intervals provide a compact statistical view of the repeated-run stability, while ECE and Brier score quantify probability calibration after the validation-derived temperature-scaling step.

\begin{table}[t]
    \centering
    \small
    \caption{Uncertainty and calibration summary for UL-TransXNet on the main test sets. Values follow the manuscript seed/fold-level reporting protocol; intervals are bootstrap 95\% confidence intervals over seeds or folds. ECE uses 10 confidence bins.}
    \label{tab:calibration_uncertainty}
    \begin{adjustbox}{max width=\columnwidth}
    \begin{tabular}{lcccc}
        \toprule
        Dataset & AUC 95\% CI & BalAcc 95\% CI & ECE & Brier \\
        \midrule
        TN5000 & 0.9483 [0.9475, 0.9488] & 0.8726 [0.8662, 0.8771] & 0.0220 & 0.0741 \\
        BUSI & 0.8998 [0.8880, 0.9067] & 0.7961 [0.7865, 0.8045] & 0.0696 & 0.1358 \\
        AUL & 0.8767 [0.8759, 0.8776] & 0.8200 [0.7755, 0.8599] & 0.1283 & 0.1597 \\
        \bottomrule
    \end{tabular}
    \end{adjustbox}
\end{table}

\subsection{Efficiency and mobile prototype}
The Android experiment evaluates a TN5000-calibrated prototype of UL-TransXNet. Table~\ref{tab:mobile_results} and Figure~\ref{fig:mobile_tradeoff} show a trade-off among three inference variants on one device. A single checkpoint provides the lowest end-to-end latency at 116.5 ms per image. Adding horizontal-flip TTA improves both AUC and balanced accuracy, reaching 0.9466 AUC and 0.8762 balanced accuracy at 226.2 ms per image. The three-checkpoint ensemble further increases AUC to 0.9485, but it reduces balanced accuracy at the selected operating threshold and raises latency to 738.6 ms per image. Thus, the single and TTA variants are the practical interactive settings for this prototype, whereas the ensemble is better interpreted as an offline ranking or review option. These measurements are not a full mobile-baseline benchmark.

\begin{table}[t]
    \centering
    \small
    \caption{Android prototype results on the TN5000 test set. Latency includes ROI preprocessing and ONNX Runtime inference.}
    \label{tab:mobile_results}
    \begin{adjustbox}{max width=\columnwidth}
    \begin{tabular}{lccccc}
        \toprule
        Mode & Runs & AUC & BalAcc & Avg ms & P95 ms \\
        \midrule
        Single & 4 & 0.9402 & 0.8688 & 116.5 & 129.0 \\
        + HFlip TTA & 2 & 0.9466 & \textbf{0.8762} & 226.2 & 246.6 \\
        3-ckpt ensemble + TTA & 2 & \textbf{0.9485} & 0.8723 & 738.6 & 763.5 \\
        \bottomrule
    \end{tabular}
    \end{adjustbox}
\end{table}

\begin{figure}[t]
    \centering
    \includegraphics[width=\textwidth]{fig_android_deployment_tradeoff.png}
    \caption{Accuracy--latency trade-off of the Android prototype variants on the TN5000 test set.}
    \label{fig:mobile_tradeoff}
\end{figure}
